\section{Visual-goal control in the deterministic gray-noise world}
\label{app:redcube-goals}

All three original \texttt{gray\_frame\_still} grid models use the same 128 freshly generated cases
and unchanged planning settings. Evaluation seed 20260920 produces native resets followed by seven
random actions; a history is admitted only if it survives this past prefix. No candidate is rejected
according to a future model result. The goal is sampled uniformly from safe positions on the same
three-pixel action lattice, with shortest Manhattan action distance between six and 12. Mean initial
L1 distance is $26.883$ pixels, with range $18$--$36$. The goal image has independently drawn gray
noise and remains fixed within its case.

The planner sees genuine eight-frame histories and the RGB goal. Horizon-three MPC enumerates all
64 native action sequences and minimizes summed squared latent-goal distance over the three
predictions. It executes the first action, observes the true successor, and replans, for at most
24 real steps. Horizon-one MPC is a predeclared comparison. No decoder, position probe, reward
head, altered reward, latent rescaling, or privileged state enters either learned planner. History
and goal encoding use a common fixed $[16,8,63,63,3]$ batch shape. Outgoing action alignment is
$[a_0,\ldots,a_6,a_{\mathrm{candidate}}]$.

The primary metric is a first exact goal hit before native death. Execution stops on the first hit
or death, and all 128 cases remain in the denominator. Secondary metrics are minimum/final L1
distance, entry within three pixels, deaths, and first-hit time. Random and blind-cycle references
share execution noise with each model. The privileged exact-transition reference searches the same
horizon with physical squared-distance costs and rejects plans with intermediate death. It reaches
every goal in the shortest number of actions; this is a protocol check, not a learned competitor.

Horizon-three exact-hit counts are $5$, $5$, and $9$ for models 0, 1, and 2. Their minimum-distance
means are $22.242$, $22.945$, and $22.031$ pixels, against $19.406$ for random actions. Native-death
counts are $9$, $8$, and $24$, against $20$ for random. Horizon-one hits no goal for any model,
with minimum-distance means $23.906$, $24.797$, and $24.047$ pixels and death counts $4$, $6$, and
$5$. The fixed cycle survives but reaches no goal. These outcomes establish failure of this
short-horizon latent-cost recipe on the specified distant-goal distribution; they are not a
positive planning control.

Independent reconstruction checks all 1,024 prefix frames per model, every action and native
transition, reachability, exact goal draws, first-hit/death absorption, summaries, visual episode
selection, and recorded hashes. Across the three fits, all cohort arrays and reference trajectories
are bit-identical. The three model evaluations do not create 384 independent cases. Visual episodes
0, 1, and 2 were fixed before results, and decoded selected forecasts are marked invalid after
terminal padding. The empty arena permits reactive goal movement; the task does not demonstrate
a need for deep planning.

\begin{table}[ht]
\centering
\caption{\textbf{The original distant-goal test remains a weak result.} Initial goals are six
to 12 native actions away. The same 128 goals and histories are used for all three models and references.
Distances are mean L1 pixels over all cases; lower is better. Random, blind-cycle, and privileged
reference trajectories are shared across model evaluations.}
\label{tab:redcube-goals}
\inputtable{redcube_goals.tex}
\end{table}

\subsection{Separately declared distance-response follow-up}
\label{app:redcube-distance}
The distance-response experiment was declared after the failed 6--12-action primary result. It
characterizes the same controller's operating range and does not replace that primary. All six
exact goal-distance shells---one, two, three, six, nine, and 12 native actions---use the same 128
original valid histories and all three existing model seeds. The H3 and H1 planners, 24-action
budget, random/blind-cycle/privileged references, native dynamics, and encoder/predictor execution
shapes remain fixed. No planner or representation parameter changes.

A newly fixed goal seed, 20260921, samples uniformly among safe reachable positions on each exact
Manhattan shell. The RGB goal background is shared across the six goals for a prefix, so only
object position changes. Execution noise and random actions are paired by original candidate ID.
The outcome is exact first hit before death. Every model--distance cell retains its own 128-case
count and Wilson interval; the 768 correlated rows are never pooled into an inferential interval.
The within-three-pixel secondary is already satisfied initially on the one-action shell and is
not evidence of movement there.

Case 0 at every shell is the predeclared visualization. Its full recorded native trajectory is
re-rendered with the recorded actions and execution keys; padding freezes after first hit or death.
The initial H12 CEM alternative was assessed but not launched. Thus any distance trend concerns
the unchanged H3 controller over different requested goals, not a tuned-horizon comparison.


All three models and the final independent audit completed. The H3 hit counts at distances
$(1,2,3,6,9,12)$ are $(128,128,121,25,3,0)$, $(128,128,123,32,2,0)$, and
$(128,128,124,34,1,0)$ for seeds 0, 1, and 2, respectively. Table~\ref{tab:goal-distance-counts}
keeps the H1 comparison and paired references alongside these outcomes. The independent audit
checks 8,179 conditions and reconstructs 176,334 executed native transitions from raw artifacts.

\begin{table}[ht]
\centering
\caption{Distance-response exact-hit counts. Initial distance is measured in native three-pixel
cardinal moves. Every cell has denominator 128. The same six shells,
prefixes and controller settings are retained for all models; no pooled binomial interval is used.}
\label{tab:goal-distance-counts}
\inputtable{redcube_goal_distance.tex}
\end{table}


\begin{figure}[p]
\centering
\paperfig{goal_distance_paths_m0.pdf}{1.0}
\caption{\textbf{Actual paths across the declared goal-distance shells.} Model seed 0, fixed case 0
in every shell. Curves are executed simulator positions, annotated only for evaluation; the planner
receives RGB history and an RGB goal. Paths stop at the first exact hit, native death, or 24 actions.
The examples and all six shells were fixed before the distance-response outcomes.}
\label{fig:goal-distance-paths}
\end{figure}

\begin{figure}[p]
\centering
\paperfig{goal_distance_rgb_m0.pdf}{1.0}
\caption{Native-color RGB views of the same preselected goal-distance executions for model 0.
These are exact keyed re-renderings of recorded real trajectories, with no model forecast or
invented continuation. The full artifact includes the corresponding displays for models 1 and 2.}
\label{fig:goal-distance-rgb}
\end{figure}
